% Chapter Template

\section{Control Strategy} % Main chapter title

%\label{Chapter9} % Change X to a consecutive number; for referencing this chapter elsewhere, use \ref{ChapterX}\ref{ChapterX}\ref{ChapterX}\ref{ChapterX}

%----------------------------------------------------------------------------------------
%	SECTION 1
%----------------------------------------------------------------------------------------
\subsection{Digital Signal Logic}
A critical aspect of ensuring precise control of the end-effector tip is the ability to accurately interpret user intent while rejecting extraneous inputs, such as physiological tremor or unintended vibrations. To achieve this, an advanced signal processing and control architecture has been implemented.

An Extended Kalman Filter (EKF) is employed to fuse and filter raw data acquired from multiple sources, including the AS5600 encoders and the feedback from the STS3009 servos. The EKF provides robust state estimation by dynamically weighting sensor inputs and predicting true user motion in the presence of noise and transient disturbances. This results in smooth, stable, and intentional control signals that faithfully represent operator commands while effectively suppressing high-frequency noise and involuntary motions.

At the core of the control system is an ESP32 microcontroller, which serves as the central computation unit. This chip is responsible for executing the EKF algorithm in real-time, processing sensory data, generating corresponding motor commands, and managing communication between system components. Its high processing power and support for floating-point operations make it well-suited for handling computationally intensive filtering and control tasks onboard.

After the user's input signals are filtered and interpreted, the refined commands are transmitted to the output-stage STS3032 motors that drive the cables. To further enhance motion accuracy and disturbance rejection, a PID control strategy is implemented based on real-time feedback from the STS3032 motors' built-in encoders. This closed-loop architecture enables precise regulation of cable tension and position, providing haptic transparency and ensuring safe, responsive operation during delicate surgical tasks. The combination of high-level EKF-based intention recognition and low-level PID motor control creates a robust hierarchical control system capable of compensating for both environmental uncertainties and human factors.

\subsection{Human-Machine Interface}

The control signals are transmitted from the rear input motors to the front output motors, ultimately driving the tip of the instrument. However, the system operates without physical switches or buttons to initiate or terminate control. To address this, a specially designed sequence of actions is implemented to ensure safe and intuitive operation.

Since the smart servos are equipped with absolute magnetic encoders, and their rotational range is limited to less than 360 degrees, the initial position of each motor can always be reliably identified. Upon powering on the device, the four front motors automatically execute a homing sequence. This process ensures that all cables are brought to a known tension state before any user input is accepted, thereby preventing uncontrolled movements and ensuring system readiness.

User control is enabled only after the homing process is complete. Slight pressure applied to the input lever serves as the trigger to activate the control linkage. This intentional design requires a deliberate action from the user, minimizing the risk of accidental engagement.

To terminate control, the user simply releases the lever. This action locks the rear input motors, effectively decoupling the user's hand from the instrument tip. This not only enhances safety by preventing unintended motions but also allows the entire instrument to be repositioned as a rigid unit during surgical procedures without affecting the end-effector's configuration. This approach provides a seamless and secure method for transitioning between active manipulation and passive positioning.